% =====================================================================
% QUESTION 5 — class hierarchy & UML diagram  (6 + 4 = 10 marks)
% Demonstrates: \umlclass, inheritance arrows, \expectedoutput
% =====================================================================
\question{5}{6 + 4 = 10}

\noindent The UML diagram below shows a three-level inheritance hierarchy
for a small shop. Implement the hierarchy in Java and complete the tasks
that follow. Coordinates given to \texttt{\textbackslash umlclass} are in
cm, measured from the origin of the diagram.

\begin{center}
\begin{tikzpicture}
    \umlclass{Product}
             {\texttt{- name : String}\\ \texttt{- price : double}}
             {\texttt{+ Product(name, price)}\\ \texttt{+ getPrice() : double}\\ \texttt{+ describe() : String}}
             {0}{0}
    \umlclass{Category}
             {\texttt{- categoryName : String}}
             {\texttt{+ Category(name, price, categoryName)}\\ \texttt{+ describe() : String}}
             {0}{-4.5}
    \umlclass{Subcategory1}
             {\texttt{- subName : String}}
             {\texttt{+ Subcategory1(name, price, categoryName, subName)}\\ \texttt{+ describe() : String}}
             {-3.5}{-9}
    \umlclass{Subcategory2}
             {\texttt{- subName : String}}
             {\texttt{+ Subcategory2(name, price, categoryName, subName)}\\ \texttt{+ describe() : String}}
             {3.5}{-9}
    \draw[inherits] (Category) -- (Product);
    \draw[inherits] (Subcategory1) -- (Category);
    \draw[inherits] (Subcategory2) -- (Category);
\end{tikzpicture}
\end{center}

\subquestion{a}{6}
Implement \texttt{Product}, \texttt{Category}, \texttt{Subcategory1} and
\texttt{Subcategory2} exactly as shown in the diagram. Every subclass must
call its parent constructor with \texttt{super(...)} and override
\texttt{describe()} so that it appends its own field to the parent's
description.

\subquestion{b}{4}
Write a \texttt{main} method that creates one object of each class, stores
all four in a \texttt{Product[]} array, and prints the result of
\texttt{describe()} for each element. The expected console output is shown
below.

\begin{expectedoutput}
> Pen (price: 5.0)\\
> Pen (price: 5.0) | Category: Stationery\\
> Pen (price: 5.0) | Category: Stationery | Sub: Writing\\
> Pen (price: 5.0) | Category: Stationery | Sub: Packing
\end{expectedoutput}
